\lecture{4}{Rotman §2.5}{Homomorphisms}

\section{Homomorphisms}

An important problem is deciding whether two given groups \(G\) and \(H\) are somehow the
same. For example, \(S_3\) is the group of all permutations of \(X = \{1,2,3\}\), and the
group \(S_Y\) of all permutations of \(Y = \{a,b,c\}\) is a different group, because
permutations of \(\{1,2,3\}\) are not permutations of \(\{a,b,c\}\). But even though
\(S_3\) and \(S_Y\) are different, they surely bear a strong resemblance to each other
(\autoref{eg:S3-SY}). The notions of homomorphism and isomorphism allow one to compare
different groups.

\subsection{Homomorphisms and Isomorphisms}

\begin{definition}[Homomorphism, isomorphism]\label{def:homomorphism}
	If \((G, *)\) and \((H, \circ)\) are groups (we have displayed the operation in each),
	then a function \(f\colon G \to H\) is a \vocab{homomorphism} if
	\[
		f(x * y) = f(x) \circ f(y)
	\]
	for all \(x, y \in G\). If \(f\) is also a bijection, then \(f\) is called an
	\vocab{isomorphism}. We say that \(G\) and \(H\) are \vocab{isomorphic}, denoted by
	\(G \cong H\), if there exists an isomorphism \(f\colon G \to H\).
\end{definition}

\begin{note}[Etymology]
	\emph{Homomorphism} comes from the Greek \emph{homo} (``same'') and \emph{morph}
	(``shape'' or ``form''): a homomorphism carries a group to another group (its image) of
	similar form. \emph{Isomorphism} involves the Greek \emph{iso} (``equal''), and
	isomorphic groups have identical form.
\end{note}

\begin{lemma}[Rotman, Exercise 2.57]\label{lem:iso-equivalence}
	\begin{enumerate}[(i)]
		\item The composite of homomorphisms is a homomorphism.
		\item The inverse of an isomorphism is an isomorphism.
		\item Isomorphism is an equivalence relation on any family of groups. In particular,
		      if \(G \cong H\), then \(H \cong G\).
	\end{enumerate}
\end{lemma}

\begin{proof}
	(i) If \(f\colon G \to H\) and \(g\colon H \to L\) are homomorphisms, then
	\(gf(xy) = g(f(x)f(y)) = gf(x)\,gf(y)\).

	(ii) Let \(f\colon G \to H\) be an isomorphism with inverse function
	\(f^{-1}\colon H \to G\). Given \(u, v \in H\), write \(u = f(x)\) and \(v = f(y)\);
	then \(f^{-1}(uv) = f^{-1}(f(x)f(y)) = f^{-1}(f(xy)) = xy = f^{-1}(u)f^{-1}(v)\).

	(iii) \(1_G\colon G \to G\) is an isomorphism (reflexivity); (ii) gives symmetry; and
	(i), together with the fact that a composite of bijections is a bijection, gives
	transitivity.
\end{proof}

\begin{eg}[First examples]\label{eg:hom-first}
	\begin{enumerate}[(i)]
		\item The identity \(1_G\colon G \to G\) is an isomorphism, and the \vocab{trivial
			      homomorphism} \(f\colon G \to H\), \(f(a) = 1\) for all \(a \in G\), is a
		      homomorphism.
		\item Let \(\mathbb R\) be the additive group of real numbers and
		      \(\mathbb R^{>}\) the multiplicative group of positive reals. Then
		      \(f\colon \mathbb R \to \mathbb R^{>}\), \(f(x) = e^x\), is a homomorphism,
		      for
		      \[
			      f(x+y) = e^{x+y} = e^xe^y = f(x)f(y).
		      \]
		      It is an isomorphism, since its inverse function is
		      \(g(x) = \log x\). Therefore \(\mathbb R \cong \mathbb R^{>}\). Note that
		      \(g\) is also an isomorphism, as \autoref{lem:iso-equivalence}(ii) predicts:
		      \(g(xy) = \log(xy) = \log x + \log y = g(x) + g(y)\).
		\item The additive group \(\mathbb C\) is isomorphic to the plane
		      \(\mathbb R^2\) (\autoref{eg:group-list}(vii)). Define
		      \(f\colon \mathbb C \to \mathbb R^2\) by \(f(a + ib) = (a,b)\). Clearly \(f\)
		      is a bijection, and it is a homomorphism:
		      \begin{align*}
			      f\bigl((a+ib) + (a'+ib')\bigr)
			       & = f\bigl((a+a') + i(b+b')\bigr) = (a+a', b+b') \\
			       & = (a,b) + (a',b') = f(a+ib) + f(a'+ib').
		      \end{align*}
	\end{enumerate}
\end{eg}

\begin{definition}[Multiplication table]
	Let \(a_1, a_2, \dots, a_n\) be a list, with no repetitions, of all the elements of a
	finite group \(G\) of order \(n\). A \vocab{multiplication table} for \(G\) is the
	\(n \times n\) array whose \(ij\) entry is \(a_ia_j\):
	\[
		\begin{array}{c|ccccc}
			G      & a_1      & \cdots & a_j      & \cdots & a_n      \\ \hline
			a_1    & a_1a_1   & \cdots & a_1a_j   & \cdots & a_1a_n   \\
			\vdots & \vdots   &        & \vdots   &        & \vdots   \\
			a_i    & a_ia_1   & \cdots & a_ia_j   & \cdots & a_ia_n   \\
			\vdots & \vdots   &        & \vdots   &        & \vdots   \\
			a_n    & a_na_1   & \cdots & a_na_j   & \cdots & a_na_n
		\end{array}
	\]
\end{definition}

We agree that the identity is listed first, \(a_1 = 1\); then the first row and first
column merely repeat the listing, and we usually omit them.

Consider two almost trivial groups: the multiplicative group \(\Gamma_2 = \{1, -1\}\) and
the parity group \(\mathbb P\) (\autoref{eg:group-list}(viii)). Their multiplication
tables are
\[
	\Gamma_2\colon\
	\begin{array}{|c|c|}
		\hline
		1  & -1 \\ \hline
		-1 & 1  \\ \hline
	\end{array}\,;
	\qquad\qquad
	\mathbb P\colon\
	\begin{array}{|c|c|}
		\hline
		\text{even} & \text{odd}  \\ \hline
		\text{odd}  & \text{even} \\ \hline
	\end{array}\,.
\]
Clearly \(\Gamma_2\) and \(\mathbb P\) are distinct groups; equally clearly, there is no
significant difference between them. Isomorphism formalizes this: the function
\(f\colon \Gamma_2 \to \mathbb P\) with \(f(1) = \text{even}\) and
\(f(-1) = \text{odd}\) is an isomorphism.

There are many multiplication tables for a group of order \(n\), one for each of the
\(n!\) lists of its elements. If \(a_1, \dots, a_n\) lists the elements of \(G\) and
\(f\colon G \to H\) is a bijection, then \(f(a_1), \dots, f(a_n)\) lists the elements of
\(H\) and determines a multiplication table for \(H\). That \(f\) is an isomorphism says
precisely that the two tables match: if \(a_ia_j\) is the \(ij\) entry of the table of
\(G\), then \(f(a_i)f(a_j) = f(a_ia_j)\) is the \(ij\) entry of the table of \(H\). In
this sense isomorphic groups have the same multiplication table; they are essentially the
same, differing only in the notation for the elements and the operations.

\begin{eg}[\(S_3\) and \(S_Y\)]\label{eg:S3-SY}
	Here is an algorithm to check whether a bijection \(f\colon G \to H\) between finite
	groups is an isomorphism: list the elements \(a_1, \dots, a_n\) of \(G\), form the
	multiplication table of \(G\) from this list and that of \(H\) from
	\(f(a_1), \dots, f(a_n)\), and compare the \(n^2\) entries.

	Take \(G = S_3\) and \(H = S_Y\), \(Y = \{a,b,c\}\). List \(S_3\) as
	\[
		(1),\ (1\ 2),\ (1\ 3),\ (2\ 3),\ (1\ 2\ 3),\ (1\ 3\ 2),
	\]
	and let \(\varphi\colon S_3 \to S_Y\) be the obvious function replacing numbers by
	letters:
	\[
		(1),\ (a\ b),\ (a\ c),\ (b\ c),\ (a\ b\ c),\ (a\ c\ b).
	\]
	Superimposing the two tables, they match. For instance, the \(4,5\) entry for \(S_3\) is
	\((2\ 3)(1\ 2\ 3) = (1\ 3)\), while the \(4,5\) entry for \(S_Y\) is
	\((b\ c)(a\ b\ c) = (a\ c)\). This is the special case \(X = \{1,2,3\}\) of
	\autoref{lem:SX-iso}.
\end{eg}

We now turn from isomorphisms to general homomorphisms.

\begin{lemma}\label{lem:hom-basics}
	Let \(f\colon G \to H\) be a homomorphism.
	\begin{enumerate}[(i)]
		\item \(f(1) = 1\);
		\item \(f(x^{-1}) = f(x)^{-1}\);
		\item \(f(x^n) = f(x)^n\) for all \(n \in \mathbb Z\).
	\end{enumerate}
\end{lemma}

\begin{proof}
	(i) Applying \(f\) to \(1 \cdot 1 = 1\) gives \(f(1)f(1) = f(1)\), and multiplying both
	sides by \(f(1)^{-1}\) gives \(f(1) = 1\).

	(ii) Applying \(f\) to \(x^{-1}x = 1\) gives \(f(x^{-1})f(x) = 1\), and uniqueness of
	inverses (\autoref{prop:group-basics}(iv)) gives \(f(x^{-1}) = f(x)^{-1}\).

	(iii) An easy induction gives \(f(x^n) = f(x)^n\) for all \(n \geq 0\). For negative
	exponents, \((y^{-1})^n = y^{-n}\) in any group, so
	\[
		f(x^{-n}) = f\bigl((x^{-1})^n\bigr) = f(x^{-1})^n = \bigl(f(x)^{-1}\bigr)^n
		= f(x)^{-n}. \qedhere
	\]
\end{proof}

\begin{eg}[Cyclic groups of the same order]\label{eg:cyclic-iso}
	Any two finite cyclic groups \(G\) and \(H\) of the same order \(m\) are isomorphic.
	Together with \autoref{cor:prime-order-cyclic}, it follows that any two groups of prime
	order \(p\) are isomorphic.

	Suppose \(G = \langle x\rangle\) and \(H = \langle y\rangle\), and define
	\(f\colon G \to H\) by \(f(x^i) = y^i\) for \(0 \leq i < m\). Since
	\(G = \{1, x, \dots, x^{m-1}\}\) and \(H = \{1, y, \dots, y^{m-1}\}\)
	(\autoref{prop:order-equals-size}), \(f\) is a bijection. To see that \(f\) is a
	homomorphism, we must show \(f(x^ix^j) = f(x^i)f(x^j)\) for \(0 \leq i, j < m\). If
	\(i + j < m\), then
	\[
		f(x^ix^j) = f(x^{i+j}) = y^{i+j} = y^iy^j = f(x^i)f(x^j).
	\]
	If \(i + j \geq m\), then \(i + j = m + r\) with \(0 \leq r < m\), so
	\(x^{i+j} = x^mx^r = x^r\) (because \(x^m = 1\)); similarly \(y^{i+j} = y^r\). Hence
	\[
		f(x^ix^j) = f(x^{i+j}) = f(x^r) = y^r = y^{i+j} = y^iy^j = f(x^i)f(x^j).
	\]
	Therefore \(f\) is an isomorphism and \(G \cong H\). (\autoref{eg:cyclic-first-iso} gives
	a nicer proof.)
\end{eg}

\begin{definition}[Invariant]
	A property of a group \(G\) that is shared by every group isomorphic to \(G\) is called an
	\vocab{invariant} of \(G\).
\end{definition}

For example, the order \(\abs{G}\) is an invariant, for isomorphic groups have the same
order. Being abelian is an invariant: if \(ab = ba\), then
\[
	f(a)f(b) = f(ab) = f(ba) = f(b)f(a),
\]
so \(f(a)\) and \(f(b)\) commute. Thus \(\mathbb R\) and \(GL(2,\mathbb R)\) are not
isomorphic, for \(\mathbb R\) is abelian and \(GL(2,\mathbb R)\) is not. Here are some
further invariants.

\begin{lemma}[Rotman, Exercise 2.59]\label{lem:order-invariant}
	Let \(f\colon G \to H\) be an isomorphism.
	\begin{enumerate}[(i)]
		\item For every \(a \in G\), the elements \(a\) and \(f(a)\) have the same order
		      (both infinite, or both finite and equal). Consequently, if \(G\) has an
		      element of some order \(n\) and \(H\) does not, then \(G \not\cong H\).
		\item For every \(k\), \(G\) and \(H\) have the same number of elements of
		      order \(k\).
	\end{enumerate}
\end{lemma}

\begin{proof}
	(i) By \autoref{lem:hom-basics}, \(f(a)^n = f(a^n)\), and since \(f\) is injective
	with \(f(1) = 1\), we have \(f(a^n) = 1\) if and only if \(a^n = 1\). Thus
	\(\{n : a^n = 1\} = \{n : f(a)^n = 1\}\), and the orders agree.

	(ii) By (i), \(f\) restricts to a bijection from the elements of order \(k\) in \(G\) to
	those in \(H\).
\end{proof}

Whether or not a group is cyclic is also an invariant. In general, however, it is a
challenge to decide whether two given groups are isomorphic.

\begin{eg}[\(\mathbf V \not\cong \Gamma_4\)]\label{eg:V-not-Gamma4}
	Here are two nonisomorphic groups of the same order. Let
	\(\mathbf V = \{(1), (1\ 2)(3\ 4), (1\ 3)(2\ 4), (1\ 4)(2\ 3)\}\) be the four-group
	(\autoref{eg:subgroups}), and let \(\Gamma_4 = \langle i\rangle = \{1, i, -1, -i\}\) be
	the multiplicative cyclic group of fourth roots of unity. If there were an isomorphism
	\(f\colon \mathbf V \to \Gamma_4\), then surjectivity would provide \(x \in \mathbf V\)
	with \(i = f(x)\). But \(x^2 = (1)\) for all \(x \in \mathbf V\), so
	\(i^2 = f(x)^2 = f(x^2) = f((1)) = 1\), contradicting \(i^2 = -1\). Therefore
	\(\mathbf V \not\cong \Gamma_4\).

	There are other proofs: \(\Gamma_4\) is cyclic and \(\mathbf V\) is not; \(\Gamma_4\)
	has an element of order \(4\) and \(\mathbf V\) does not; \(\Gamma_4\) has a unique
	element of order \(2\), while \(\mathbf V\) has three.
\end{eg}

\subsection{Kernels, Images, and Normal Subgroups}

\begin{definition}[Kernel, image]
	If \(f\colon G \to H\) is a homomorphism, define its \vocab{kernel} and \vocab{image} by
	\[
		\ker f = \{x \in G : f(x) = 1\}
		\qquad\text{and}\qquad
		\im f = \{h \in H : h = f(x) \text{ for some } x \in G\}.
	\]
\end{definition}

\begin{note}
	\emph{Kernel} comes from the German word for ``grain'' or ``seed'' (\emph{corn} comes
	from the same word): it indicates an important ingredient of a homomorphism.
\end{note}

\begin{eg}[Kernels and images]\label{eg:kernels}
	\begin{enumerate}[(i)]
		\item Let \(\Gamma_n = \langle\zeta\rangle\), where \(\zeta = e^{2\pi i/n}\) is a
		      primitive \(n\)th root of unity. Then \(f\colon \mathbb Z \to \Gamma_n\),
		      \(f(m) = \zeta^m\), is a surjective homomorphism, and \(\ker f\) consists of
		      all the multiples of \(n\), because \(\zeta\) has order \(n\)
		      (\autoref{lem:order-divides}).
		\item \(\sgn\colon S_n \to \Gamma_2 = \{\pm 1\}\) is a homomorphism by
		      \autoref{thm:sgn-multiplicative}. For \(n \geq 2\) it is surjective, since
		      \(\sgn(\tau) = -1\) for a transposition \(\tau\); its kernel is the alternating
		      group \(A_n\) of all even permutations.
		\item Since \(\det(AB) = \det(A)\det(B)\), the determinant
		      \(\det\colon GL(2,\mathbb R) \to \mathbb R^\times\) is a homomorphism. It is
		      surjective, since
		      each \(r \in \mathbb R^\times\) is the determinant of
		      \(\begin{psmallmatrix} r & 0 \\ 0 & 1 \end{psmallmatrix}\). Its kernel is the
		      \vocab{special linear group}
		      \[
			      SL(2,\mathbb R) = \{A \in GL(2,\mathbb R) : \det(A) = 1\}.
		      \]
		      (This extends to \(\det\colon GL(n,\mathbb R) \to \mathbb R^\times\).)
		\item We generalize the construction of \(\ker f\). Recall the \vocab{inverse image}
		      of a subset \(B \subseteq Y\) under a function \(f\colon X \to Y\):
		      \[
			      f^{-1}(B) = \{x \in X : f(x) \in B\}.
		      \]
		      If \(f\colon G \to H\) is a homomorphism and \(B \leq H\), then
		      \(f^{-1}(B)\) is a subgroup of \(G\) containing \(\ker f\). Indeed,
		      \(1 \in f^{-1}(B)\), for \(f(1) = 1 \in B\). If \(x, y \in f^{-1}(B)\), then
		      \(f(xy) = f(x)f(y) \in B\), so \(xy \in f^{-1}(B)\). If \(x \in f^{-1}(B)\),
		      then \(f(x^{-1}) = f(x)^{-1} \in B\), so \(x^{-1} \in f^{-1}(B)\). Finally
		      \(\ker f = f^{-1}(\{1\}) \subseteq f^{-1}(B)\).
	\end{enumerate}
\end{eg}

\begin{proposition}\label{prop:kernel-image}
	Let \(f\colon G \to H\) be a homomorphism.
	\begin{enumerate}[(i)]
		\item \(\ker f\) is a subgroup of \(G\), and \(\im f\) is a subgroup of \(H\).
		\item If \(x \in \ker f\) and \(a \in G\), then \(axa^{-1} \in \ker f\).
		\item \(f\) is an injection if and only if \(\ker f = \{1\}\).
	\end{enumerate}
\end{proposition}

\begin{proof}
	(i) \(\ker f = f^{-1}(\{1\})\) is a subgroup by \autoref{eg:kernels}(iv). As for the
	image: \(1 = f(1) \in \im f\); if \(h = f(x) \in \im f\), then
	\(h^{-1} = f(x)^{-1} = f(x^{-1}) \in \im f\); and if also \(k = f(y) \in \im f\), then
	\(hk = f(x)f(y) = f(xy) \in \im f\).

	(ii) If \(f(x) = 1\), then
	\(f(axa^{-1}) = f(a)f(x)f(a)^{-1} = f(a)f(a)^{-1} = 1\).

	(iii) If \(f\) is an injection, then \(x \neq 1\) implies \(f(x) \neq f(1) = 1\), so
	\(x \notin \ker f\). Conversely, assume \(\ker f = \{1\}\) and \(f(x) = f(y)\). Then
	\(1 = f(x)f(y)^{-1} = f(xy^{-1})\), so \(xy^{-1} \in \ker f = \{1\}\); hence
	\(x = y\).
\end{proof}

\begin{definition}[Normal subgroup]\label{def:normal}
	A subgroup \(K\) of a group \(G\) is a \vocab{normal subgroup} if \(k \in K\) and
	\(g \in G\) imply \(gkg^{-1} \in K\). We then write \(K \trianglelefteq G\).
\end{definition}

\autoref{prop:kernel-image} thus says that \emph{the kernel of a homomorphism is always a
	normal subgroup}. We record some first examples.
\begin{itemize}
	\item If \(G\) is abelian, then every subgroup \(K\) is normal: for \(k \in K\) and
	      \(g \in G\), \(gkg^{-1} = kgg^{-1} = k \in K\).
	\item The cyclic subgroup \(H = \langle(1\ 2)\rangle = \{(1), (1\ 2)\}\) of \(S_3\) is
	      \emph{not} normal: if \(\alpha = (1\ 2\ 3)\), then \(\alpha^{-1} = (3\ 2\ 1)\) and
	      \[
		      \alpha(1\ 2)\alpha^{-1} = (1\ 2\ 3)(1\ 2)(3\ 2\ 1) = (2\ 3) \notin H.
	      \]
	\item The cyclic subgroup \(K = \langle(1\ 2\ 3)\rangle = \{(1), (1\ 2\ 3), (1\ 3\ 2)\}\)
	      of \(S_3\) is normal: by \autoref{prop:conj-same-structure}, a conjugate of a
	      \(3\)-cycle is a \(3\)-cycle, and \(K\) contains both \(3\)-cycles of \(S_3\).
	\item By \autoref{eg:kernels}(ii) and (iii), \(A_n \trianglelefteq S_n\) and
	      \(SL(2,\mathbb R) \trianglelefteq GL(2,\mathbb R)\) (both facts are also easy to
	      prove directly).
\end{itemize}

\subsection{Conjugation and the Center}

\begin{definition}[Conjugate]
	If \(G\) is a group and \(a \in G\), then a \vocab{conjugate} of \(a\) is any element of
	the form \(gag^{-1}\) with \(g \in G\).
\end{definition}

Clearly, a subgroup \(K \leq G\) is normal if and only if \(K\) contains all the conjugates
of its elements. By \autoref{prop:conjugacy-criterion}, \(\alpha, \beta \in S_n\) are
conjugate in \(S_n\) if and only if they have the same cycle structure.

\begin{remark}
	If \(H \leq S_n\), then \(\alpha, \beta \in H\) being conjugate in \(S_n\) does
	\emph{not} imply that they are conjugate in \(H\). For example, \((1\ 2)(3\ 4)\) and
	\((1\ 3)(2\ 4)\) are conjugate in \(S_4\), but not in \(\mathbf V\): since
	\(\mathbf V\) is abelian, the only conjugate of an element of \(\mathbf V\) in
	\(\mathbf V\) is the element itself.
\end{remark}

\begin{remark}[Similarity]
	In linear algebra, a linear transformation \(T\colon V \to V\) of an \(n\)-dimensional
	real vector space determines an \(n \times n\) matrix \(A\) once a basis of \(V\) is
	chosen; another basis gives another matrix \(B\), and \(A\) and \(B\) are
	\emph{similar}: \(B = PAP^{-1}\) for some nonsingular \(P\). Thus conjugacy in
	\(GL(n,\mathbb R)\) is similarity.
\end{remark}

\begin{definition}[Conjugation]
	If \(G\) is a group and \(g \in G\), define \vocab{conjugation}
	\(\gamma_g\colon G \to G\) by
	\[
		\gamma_g(a) = gag^{-1} \qquad\text{for all } a \in G.
	\]
\end{definition}

\begin{proposition}\label{prop:conjugation-iso}
	\begin{enumerate}[(i)]
		\item If \(G\) is a group and \(g \in G\), then conjugation
		      \(\gamma_g\colon G \to G\) is an isomorphism.
		\item Conjugate elements have the same order.
	\end{enumerate}
\end{proposition}

\begin{proof}
	(i) If \(g, h \in G\), then
	\[
		(\gamma_g \circ \gamma_h)(a) = \gamma_g(hah^{-1}) = g(hah^{-1})g^{-1}
		= (gh)a(gh)^{-1} = \gamma_{gh}(a);
	\]
	that is, \(\gamma_g \circ \gamma_h = \gamma_{gh}\). It follows that each \(\gamma_g\) is
	a bijection, for \(\gamma_g \circ \gamma_{g^{-1}} = \gamma_1 = 1_G
	= \gamma_{g^{-1}} \circ \gamma_g\). And \(\gamma_g\) is a homomorphism:
	\[
		\gamma_g(ab) = g(ab)g^{-1} = (gag^{-1})(gbg^{-1}) = \gamma_g(a)\gamma_g(b).
	\]

	(ii) If \(a\) and \(b\) are conjugate, then \(b = gag^{-1} = \gamma_g(a)\) for some
	\(g \in G\). Since \(\gamma_g\) is an isomorphism, \(a\) and \(b\) have the same order by
	\autoref{lem:order-invariant}(i).
\end{proof}

\begin{eg}[The center]\label{eg:center}
	The \vocab{center} of a group \(G\) is
	\[
		Z(G) = \{z \in G : zg = gz \text{ for all } g \in G\};
	\]
	it consists of all elements commuting with everything in \(G\). (The equation
	\(zg = gz\) can be rewritten as \(z = gzg^{-1}\), so no other element of \(G\) is
	conjugate to \(z\).)

	\(Z(G)\) is a subgroup: \(1 \in Z(G)\). If \(y, z \in Z(G)\), then for all \(g\),
	\((yz)g = y(zg) = y(gz) = (yg)z = g(yz)\), so \(yz \in Z(G)\). If \(z \in Z(G)\), then
	in particular \(zg^{-1} = g^{-1}z\) for all \(g\), and taking inverses
	(\autoref{lem:group-cancellation}(iii)) gives \(gz^{-1} = z^{-1}g\), so
	\(z^{-1} \in Z(G)\).

	\(Z(G)\) is a normal subgroup: if \(z \in Z(G)\) and \(g \in G\), then
	\(gzg^{-1} = zgg^{-1} = z \in Z(G)\). The same computation shows that \emph{every}
	subgroup of \(Z(G)\) is normal in \(G\).

	A group \(G\) is abelian if and only if \(Z(G) = G\). At the other extreme are groups with
	\(Z(G) = \{1\}\); these are called \vocab{centerless}. For example, \(S_n\) is centerless
	for every \(n \geq 3\) (Rotman, Exercise 2.31): if \(\alpha \neq (1)\), then \(\alpha\)
	moves some \(i\), say \(\alpha(i) = j \neq i\); choose \(k \notin \{i, j\}\) and let
	\(\beta = (j\ k)\). By \autoref{lem:conj-action}, \(\beta\alpha\beta^{-1}\) sends
	\(\beta(i) = i\) to \(\beta(j) = k \neq j\), so \(\beta\alpha\beta^{-1} \neq \alpha\),
	and \(\alpha \notin Z(S_n)\).
\end{eg}

\begin{eg}[\(\mathbf V \trianglelefteq S_4\)]\label{eg:V-normal-S4}
	The four-group \(\mathbf V = \{(1), (1\ 2)(3\ 4), (1\ 3)(2\ 4), (1\ 4)(2\ 3)\}\) is a
	normal subgroup of \(S_4\). By \autoref{prop:conj-same-structure}, every conjugate of a
	product of two disjoint transpositions is another such; but there are only \(3\)
	permutations in \(S_4\) with this cycle structure (\autoref{tab:cycle-counts}), and all
	of them lie in \(\mathbf V\).
\end{eg}

\begin{proposition}\label{prop:index-two}
	Let \(H\) be a subgroup of index \(2\) in a group \(G\).
	\begin{enumerate}[(i)]
		\item \(g^2 \in H\) for every \(g \in G\).
		\item \(H\) is a normal subgroup of \(G\).
	\end{enumerate}
\end{proposition}

\begin{proof}
	Since \(H\) has index \(2\), there are exactly two cosets, \(H\) and \(aH\) with
	\(a \notin H\), and \(G\) is the disjoint union \(G = H \cup aH\).

	(i) If \(g \in H\), then \(g^2 \in H\). Otherwise \(g \in aH\), say \(g = ah\) with
	\(h \in H\). If \(g^2 \notin H\), then \(g^2 = ah'\) with \(h' \in H\), and
	\[
		g = g^{-1}g^2 = (ah)^{-1}ah' = h^{-1}a^{-1}ah' = h^{-1}h' \in H,
	\]
	a contradiction.

	(ii) Let \(h \in H\) and \(g \in G\); we show \(ghg^{-1} \in H\). If \(g \in H\), this
	holds because \(H\) is a subgroup. Otherwise \(g \in aH\), say \(g = ax\) with
	\(x \in H\). Then \(ghg^{-1} = a(xhx^{-1})a^{-1} = ah'a^{-1}\), where
	\(h' = xhx^{-1} \in H\). If \(ghg^{-1} \notin H\), then \(ah'a^{-1} \in aH\), say
	\(ah'a^{-1} = ay\) with \(y \in H\). Canceling \(a\) gives \(h'a^{-1} = y\), and then
	\(a = y^{-1}h' \in H\), a contradiction. Hence \(H \trianglelefteq G\).
\end{proof}

\subsection{The Quaternions}

\begin{definition}[Quaternions]\label{def:quaternions}
	The \vocab{group of quaternions} is the set \(\mathbf Q\) of the following eight matrices
	in \(GL(2,\mathbb C)\):
	\[
		\mathbf Q = \{I, A, A^2, A^3, B, BA, BA^2, BA^3\},
		\qquad\text{where}\quad
		A = \begin{pmatrix} 0 & 1 \\ -1 & 0 \end{pmatrix},
		\quad
		B = \begin{pmatrix} 0 & i \\ i & 0 \end{pmatrix}.
	\]
\end{definition}

\begin{proposition}[Rotman, Exercise 2.76]\label{prop:quaternions}
	\(\mathbf Q\) is a nonabelian group of order \(8\). Its only element of order \(2\) is
	\(-I\), every \(M \in \mathbf Q\) with \(M \neq \pm I\) satisfies \(M^2 = -I\), and
	\(Z(\mathbf Q) = \{I, -I\}\).
\end{proposition}

\begin{proof}
	Put \(C = BA\). Direct computation gives
	\[
		A^2 = B^2 = -I, \qquad
		C = BA = \begin{pmatrix} -i & 0 \\ 0 & i \end{pmatrix}, \qquad
		AB = \begin{pmatrix} i & 0 \\ 0 & -i \end{pmatrix} = -C, \qquad C^2 = -I.
	\]
	Hence \(A^3 = -A\), \(BA^2 = -B\), \(BA^3 = -C\), and
	\(\mathbf Q = \{\pm I, \pm A, \pm B, \pm C\}\), eight distinct matrices. From the
	relations above,
	\begin{align*}
		AB & = -C,             & BC & = B^2A = -A,        & CA & = BA^2 = -B,     \\
		BA & = C,              & CB & = BAB = -B^2A = A,  & AC & = ABA = -CA = B,
	\end{align*}
	Thus the product of any two elements of \(\mathbf Q\) lies in
	\(\mathbf Q\), and \(\mathbf Q\) is a subgroup of \(GL(2,\mathbb C)\) by
	\autoref{prop:finite-subgroup-criterion}. It is not abelian, since \(AB = -BA \neq BA\).

	Since \((-M)^2 = M^2\), every element other than \(\pm I\) squares to \(-I\) and hence has
	order \(4\), while \(-I\) has order \(2\). Finally \(\pm I\) are central (they are scalar
	matrices), whereas \(\pm A\) does not commute with \(B\), and \(\pm B\) and \(\pm C\) do
	not commute with \(A\); so \(Z(\mathbf Q) = \{\pm I\}\).
\end{proof}

\begin{note}[Hamilton's quaternions]
	Addition, subtraction, multiplication, and division can be extended from \(\mathbb R\) to
	the plane so that all the usual laws hold: this is \(\mathbb C\). W.~R.~Hamilton found
	a way of extending these operations to four-dimensional space so that all the usual
	laws still hold except commutativity of multiplication; he called the new ``numbers''
	\emph{quaternions} (Latin for ``four''). The multiplication is determined by the basis
	vectors \(1, \mathbf i, \mathbf j, \mathbf k\) with
	\[
		\mathbf i^2 = \mathbf j^2 = \mathbf k^2 = -1, \quad
		\mathbf{ij} = \mathbf k = -\mathbf{ji}, \quad
		\mathbf{jk} = \mathbf i = -\mathbf{kj}, \quad
		\mathbf{ki} = \mathbf j = -\mathbf{ik}.
	\]
	The nonzero quaternions form a multiplicative group, and \(\mathbf Q\) is isomorphic to
	its subgroup \(\{\pm 1, \pm\mathbf i, \pm\mathbf j, \pm\mathbf k\}\).
\end{note}

In \(\mathbf Q\), the element \(A\) has order \(4\), so \(\langle A\rangle\) has index
\(2\); its other coset is
\[
	B\langle A\rangle = \{B, BA, BA^2, BA^3\}.
\]

\begin{eg}[Every subgroup of \(\mathbf Q\) is normal]\label{eg:Q-subgroups-normal}
	By Lagrange's theorem the possible orders of subgroups of \(\mathbf Q\) are \(1, 2, 4,
	8\). The subgroups \(\{I\}\) and \(\mathbf Q\) are clearly normal. A subgroup of order
	\(4\) has index \(2\), hence is normal by \autoref{prop:index-two}. Finally, \(-I\) is
	the only element of order \(2\), so \(\langle -I\rangle\) is the only subgroup of order
	\(2\); it is normal because it equals the center \(Z(\mathbf Q)\)
	(\autoref{eg:center}).
\end{eg}

Thus \(\mathbf Q\) is a nonabelian group which is like an abelian group in the sense that
every subgroup is normal. This is essentially the only such example: every finite group
all of whose subgroups are normal has the form \(\mathbf Q \times A\), where
\(A = B \times C\), every nonidentity element of \(B\) has order \(2\), and every element
of \(C\) has odd order (direct products are introduced in the next section).

Lagrange's theorem says that the order of a subgroup of a finite group \(G\) divides
\(\abs{G}\). As remarked after \autoref{cor:prime-order-cyclic}, the converse is false; we
can now prove it.

\begin{proposition}\label{prop:A4-no-6}
	The alternating group \(A_4\) is a group of order \(12\) having no subgroup of order
	\(6\).
\end{proposition}

\begin{proof}
	By \autoref{tab:cycle-counts}, the even permutations in \(S_4\) are the identity, the
	\(8\) three-cycles, and the \(3\) products of two disjoint transpositions; so
	\(\abs{A_4} = 12\). If \(A_4\) had a subgroup \(H\) of order \(6\), then \(H\) would have
	index \(2\), and so \(\alpha^2 \in H\) for every \(\alpha \in A_4\), by
	\autoref{prop:index-two}(i). If \(\alpha\) is a \(3\)-cycle, then \(\alpha\) has order
	\(3\), so \(\alpha = \alpha^4 = (\alpha^2)^2 \in H\). Thus \(H\) contains all \(8\)
	three-cycles, which is impossible since \(\abs{H} = 6\).
\end{proof}

\autoref{prop:abelian-subgroups} will show, however, that an \emph{abelian} group of order
\(n\) has a subgroup of order \(d\) for every divisor \(d\) of \(n\).
